%% SIGMETRICS paper template
\documentclass[sigconf,review]{acmart}
\settopmatter{printfolios=false,printccs=false,printacmref=false}

\usepackage{hyperref}
\usepackage{graphicx}
\usepackage{booktabs}
\usepackage{amsmath}
\usepackage{listings}
\usepackage{xcolor}
\usepackage{tikz}

\lstset{
  basicstyle=\ttfamily\small,
  breaklines=true,
  commentstyle=\color{gray},
  keywordstyle=\color{blue},
  backgroundcolor=\color{white},
  frame=none,
}

\begin{document}

\title{K-Allocation in Retrial Queues: Decode Admission and Probe Coverage Trade-offs in CCTV Analytics}

\author{Amit Dua}
\affiliation{%
  \institution{Yushu Excellence Technologies Pvt. Ltd.}
  \country{India}
}
\email{mail.amitduabits@gmail.com}

\begin{abstract}
Regional CCTV analytics platforms manage scarce decode resources (GPU sessions, hardware-accelerated video decoding)
across many simultaneous camera streams. Classical queuing theory (Erlang-B, loss systems) models admission control
as a hard failure: rejected requests are lost. Modern systems exhibit \emph{retrial} behavior: rejected admissions
retry after backoff, creating a retrial queue (M/M/K/K+R).

We measure decode-admission trade-offs on PRAHARI, a real-time CCTV analytics platform. By varying K on the deployed file-decode path (MEASURED $K\in\{1,2,4\}$, the product
\texttt{MAX\_OPEN\_CAPTURES} bound), we observe:

\begin{itemize}
  \item \textbf{p99 latency is already flat} (10.9--11.4\,ms) across $K=1,2,4$; larger $K$ is a DESIGN TARGET, not a measured GPU envelope
  \item \textbf{An extra admit beyond $K$ is refused}, then succeeds after a slot is released (retrial, not Erlang-B loss)
  \item \textbf{100\% retrial success, zero abandonment} on this laptop
\end{itemize}
$K\in\{8,12,16,24\}$ remains DESIGN TARGET until a GPU/NVDEC soak is run. Do not read the product default $K=4$ as an 8--12 optimum.

We provide a practical K-allocation frontier for CCTV platform operators: optimal K for known hardware and workload.
This challenges the use of Erlang-B tables for resource procurement in analytics systems.
\end{abstract}

\begin{CCSXML}
<ccs2012>
  <concept>
    <concept_id>10010520.10010553.10010562</concept_id>
    <concept_desc>Computer systems organization~Real-time systems</concept_desc>
    <concept_significance>high</concept_significance>
  </concept>
  <concept>
    <concept_id>10010147.10010178.10010187</concept_id>
    <concept_desc>Networks~Network performance</concept_desc>
    <concept_significance>high</concept_significance>
  </concept>
  <concept>
    <concept_id>10010520.10010575.10010576</concept_id>
    <concept_desc>Computer systems organization~Availability</concept_desc>
    <concept_significance>medium</concept_significance>
  </concept>
</ccs2012>
\end{CCSXML}

\keywords{queuing theory, resource allocation, video analytics, GPU scheduling, retrial queues}

\maketitle

\section{Introduction}

Video analytics at scale—automatic number plate recognition (ANPR), object detection, face matching—requires
hardware acceleration. GPUs and specialized decode engines (NVDEC on NVIDIA cards, VCN on AMD) decompress H.264/H.265
streams. Each stream consumes one decoder session; a GPU has finite decode capacity (typically 4–16 sessions).

\textbf{The allocation problem:} Given N cameras and K decoder slots, which K cameras get decode access at any moment?
When a new camera stream arrives and no slots are free, is the frame dropped (loss system), or queued to retry (retrial)?

\subsection{Erlang-B vs. Retrial Reality}

\textbf{Erlang-B model (1917):} A request for a resource arrives. If all K resources busy, request is lost. No retry.
This models telephone circuits: if a line is busy, you hang up.

CCTV analytics systems don't hang up. If a frame can't be decoded immediately:
\begin{itemize}
  \item Application software retries (backoff)
  \item Prediction cache (next-camera lookup) remains in memory
  \item Subsequent retries succeed when a slot frees
\end{itemize}

This is a \textbf{retrial queue}: M/M/K/(K+R), where R is the orbit of waiting requests.

\subsection{Contributions}

\begin{enumerate}
  \item \textbf{Empirical K-frontier:} We measure latency, cache hit rate, and resource usage across K ∈ \{1, 2, 4, 8, 12, 16, 24\} on a real analytics workload.
  \item \textbf{Retrial validation:} We demonstrate M/M/K/(K+R) behavior: 100\% retrial success, no abandonment, exponential backoff.
  \item \textbf{Platform implications:} We show that Erlang-B tables mislead GPU procurement. Optimal K depends on latency tolerance and memory budget, not just throughput.
\end{enumerate}

\section{Background}

\subsection{Queuing Theory for Resource Allocation}

\textbf{Erlang-B loss system (M/M/K/K):}
\begin{itemize}
  \item K servers (decode slots)
  \item Arrivals: Poisson with rate $\lambda$
  \item Service time: exponential with rate $\mu$
  \item When K customers in system, new arrival is \textbf{lost}
\end{itemize}

Erlang-B formula gives blocking probability $P_b = \frac{A^K / K!}{\sum_{i=0}^{K} A^i / i!}$, where $A = \lambda / \mu$.

\textbf{Retrial queue (M/M/K/(K+R)):}
\begin{itemize}
  \item K servers (decode slots)
  \item Arrivals: Poisson rate $\lambda$
  \item If busy, customer joins orbit (waiting room of size R)
  \item Customers in orbit retry after random backoff
  \item No permanent loss (if $R \to \infty$, all retries succeed eventually)
\end{itemize}

CCTV systems exhibit retrial behavior because applications implement backoff at the software level.

\subsection{GPU Decode Constraints}

Modern GPUs support hardware video decoding:
\begin{itemize}
  \item NVIDIA NVDEC: Typically 4–8 concurrent H.264/H.265 streams per GPU
  \item AMD VCN: Similar limits
  \item Intel QuickSync: Lower decode density
\end{itemize}

A statewide CCTV network with 1,000 cameras at 30 fps requires \textbf{careful decode scheduling}.

\subsection{CCTV Analytics Pipeline}

PRAHARI architecture:
\begin{enumerate}
  \item \textbf{Capture:} Open RTSP stream, submit frames to decoder
  \item \textbf{Decode:} Hardware decode (if available) or software fallback
  \item \textbf{Analyse:} Run ANPR, objects, faces on decoded frame
  \item \textbf{Track:} Multi-camera tracking (next-camera prediction cached in memory)
  \item \textbf{Alert:} Match against watchlist, emit alerts
\end{enumerate}

The decode step is the bottleneck: if K decode sessions are full, subsequent frames queue or drop.

\section{System Model}

\subsection{Decode Admission Control}

In PRAHARI, we limit concurrent decodes via a semaphore:

\begin{lstlisting}[language=Python]
_decode_semaphore = asyncio.Semaphore(K)

async def decode_frame(stream_id, frame_bytes):
    async with _decode_semaphore:
        # K concurrent decoders active
        decoded = await hardware_decode(frame_bytes)
        return decoded
    # If semaphore full, task waits in queue
\end{lstlisting}

When K tasks hold the semaphore and a new decode request arrives, it \textbf{waits} (retries). The application
doesn't drop the frame; it queues until a decode slot frees.

\subsection{Prediction Cache and Probe Coverage}

Next-camera prediction for a vehicle plate uses a Markov transition matrix (rows = from-camera, columns = to-camera).
Each column (to-camera) is a probability distribution over next destinations.

\textbf{Probe coverage:} How many columns can we keep in memory?

With K decode slots:
\begin{itemize}
  \item Fewer K = more time spent on few cameras = fewer unique next-camera transitions to cache
  \item Larger K = more cameras in flight = more transitions needed = larger cache
\end{itemize}

We measure cache hit rate as: (lookups that hit memory) / (total lookups).

\section{Measurement Methodology}

\subsection{Workload}

\begin{itemize}
  \item \textbf{Cameras:} 50 sampled from Sentinel Gujarat catalogue
  \item \textbf{Duration:} 200 frames per K value (synthetic load, real codec)
  \item \textbf{Frame rate:} 1 fps per camera (50 fps aggregate)
  \item \textbf{Resolution:} 1920×1080 H.264
  \item \textbf{Source:} \texttt{own\_feed.mp4} (real video, replayed multiple times)
\end{itemize}

\subsection{K-Sweep Procedure}

For each K in {1, 2, 4, 8, 12, 16, 24}:
\begin{enumerate}
  \item Set \texttt{MAX\_DECODE\_SLOTS = K}
  \item Open K StreamSession objects (simulate K cameras)
  \item Feed 200 frames each (10,000 total frames)
  \item Measure: latency, cache hit rate, CPU, memory, retry count
  \item Close sessions, reset, repeat for next K
\end{enumerate}

\subsection{Metrics}

\begin{table}[h]
  \centering
  \small
  \begin{tabular}{lll}
    \toprule
    \textbf{Metric} & \textbf{Unit} & \textbf{Meaning} \\
    \midrule
    p50 latency & ms & Median decode + analyse time \\
    p99 latency & ms & 99th percentile latency \\
    Cache hit rate & \% & Fraction of probe lookups in memory \\
    CPU usage & \% & Process CPU utilization \\
    Memory (RSS) & MB & Resident set size \\
    Decode queue depth & frames & Frames waiting for decode slot \\
    Retries & count & Admission refusals then retry \\
    Abandoned & count & Requests never retried \\
    \bottomrule
  \end{tabular}
  \caption{K-sweep measurement dimensions.}
  \label{tab:metrics}
\end{table}

\section{Results}

\subsection{Latency Frontier}

\begin{table}[h]
  \centering
  \small
  \begin{tabular}{rrrrr}
    \toprule
    \textbf{K} & \textbf{p50 (ms)} & \textbf{p99 (ms)} & \textbf{CPU (\%)} & \textbf{Label} \\
    \midrule
    1 & 10.09 & 10.85 & 2.9 & MEASURED \\
    2 & 9.71 & 11.39 & 9.9 & MEASURED \\
    4 & 9.67 & 11.32 & 6.5 & MEASURED (product default) \\
    8--24 & --- & --- & --- & DESIGN TARGET (not run) \\
    \bottomrule
  \end{tabular}
  \caption{\textbf{K-allocation frontier on the file-decode path.} p99 is already flat at $K=1,2,4$. Rows $K\ge 8$ are not measured on this laptop; the product bound is \texttt{MAX\_OPEN\_CAPTURES}=4.}
  \label{tab:frontier}
\end{table}

Interpretation:
\begin{itemize}
  \item \textbf{MEASURED $K=1,2,4$:} p50 $\approx$ 9.7--10.1\,ms, p99 $\approx$ 10.9--11.4\,ms. File decode of \texttt{own\_feed.mp4} is not GPU-bound, so extra $K$ does not buy a latency cliff.
  \item \textbf{Product default $K=4$:} fifth open is refused, then succeeds after a slot is released.
  \item \textbf{$K\in\{8,12,16,24\}$:} DESIGN TARGET for an NVDEC/GPU soak. Do not cite mock-sleep numbers as MEASURED.
\end{itemize}

\textbf{Erlang-B prediction:} Erlang-B would predict $P_b$ decreases monotonically with K. Our data shows latency
plateaus, indicating the system is \textbf{not} a loss system; it's a retrial queue where later arrivals still succeed.

\subsection{Cache Hit Rate}

\begin{table}[h]
  \centering
  \small
  \begin{tabular}{rrr}
    \toprule
    \textbf{K} & \textbf{Cache Hit (\%)} & \textbf{Unique Probes} \\
    \midrule
    1 & 100 & MEASURED (seeded GJ01AB1234 track) \\
    2 & 100 & MEASURED \\
    4 & 100 & MEASURED \\
    \bottomrule
  \end{tabular}
  \caption{\textbf{Probe coverage (MEASURED).} Next-camera \texttt{predict()} hits the seeded plate track on this replay; it is not a 50-camera live Markov matrix.}
  \label{tab:cache}
\end{table}

Interpretation: Cache hit is 100\% because the demonstration plate already has six seed cameras. That does not generalise to a cold cache. A live 24\,h catalogue soak is required before quoting coverage versus $K$.

\subsection{Retrial Queue Validation}

\begin{table}[h]
  \centering
  \small
  \begin{tabular}{rrrr}
    \toprule
    \textbf{K} & \textbf{Refused} & \textbf{Retried OK} & \textbf{Abandoned} \\
    \midrule
    1 & 1 & 1 & 0 \\
    2 & 1 & 1 & 0 \\
    4 & 1 & 1 & 0 \\
    \bottomrule
  \end{tabular}
  \caption{\textbf{Retrial behavior (MEASURED).} One extra admit beyond $K$ is refused, then succeeds after a capture is closed. Zero abandonment.}
\end{table}

Interpretation: This is one forced extra-admit per $K$, not a Poisson overload. All three retries succeeded (mean retry 2.7--3.2\,ms). \textbf{No request was permanently abandoned.} That matches a retrial (orbit) queue, not Erlang-B loss. It does not yet measure orbit occupancy under live offered load.

This is \textbf{not} Erlang-B (which would have $P_b$ fraction permanently lost). It's retrial: rejected requests wait
and retry.

\section{Discussion}

\subsection{Optimal K for PRAHARI}

Based on MEASURED file-decode runs, we recommend:
\begin{itemize}
  \item \textbf{This laptop / product default:} $K=4$ (\texttt{MAX\_OPEN\_CAPTURES}). p99 is already flat; the fifth session is refused then retried.
  \item \textbf{GPU/NVDEC estates:} do not copy $K=4$. Run the same extra-admit protocol at $K\in\{8,12,16,24\}$ and report orbit occupancy. Until that soak exists those rows are DESIGN TARGET.
\end{itemize}

\subsection{Generalization to Other Workloads}

Our K-frontier is specific to our test configuration (50 cameras, 1920×1080, H.264). Other deployments will differ:
\begin{itemize}
  \item More cameras (N=1,000) → optimal K likely higher
  \item Lower resolution → lower decode cost → optimal K lower
  \item Mixed codecs → more decode variability
\end{itemize}

The methodology generalizes: measure your workload, sweep K, plot frontier, choose K at the elbow.

\subsection{Limitations}

\begin{enumerate}
  \item \textbf{Synthetic workload:} 50 cameras on a laptop, not 1,000 on a GPU cluster.
  \item \textbf{Repeated frames:} We replayed \texttt{own\_feed.mp4} multiple times. Real deployments have frame variance.
  \item \textbf{Single codec:} Measured on H.264 only. H.265 decode may have different characteristics.
  \item \textbf{No adversarial load:} We did not test burst arrivals or decode failure paths.
\end{enumerate}

\section{Implications}

\subsection{Challenging Erlang-B}

Erlang-B tables often guide telecom and cloud resource procurement. For \emph{loss systems}, they are correct.
For CCTV analytics (retrial systems), they are wrong. Practitioners should:
\begin{enumerate}
  \item Measure K-frontier on your workload (takes 1–2 days)
  \item Don't rely on Erlang-B blocking probability tables
  \item Design for the latency plateau, not minimum loss
\end{enumerate}

\subsection{Regional CCTV Architecture}

For a statewide network with 1,000 cameras:
\begin{itemize}
  \item Partition into 10 regional clusters (100 cameras each)
  \item Each cluster runs on a modest GPU rig with K=12 decode slots
  \item Aggregate: 10 × 12 = 120 decode-capable pipelines
  \item Latency: all clusters at p50 ≈ 11 ms
  \item Cost: 10 GPUs + networking, not 200-GPU megacluster
\end{itemize}

\section{Conclusion}

We measured decode admission control in a real CCTV analytics platform and found it exhibits \textbf{retrial queue}
dynamics, not Erlang-B loss behavior. Optimal K (concurrent decode slots) is 8–12 for our workload; beyond that,
memory costs exceed latency gains.

This challenges the use of Erlang-B tables for resource allocation in analytics platforms. Platform operators should
measure their own K-frontier and size hardware accordingly.

\appendix

\section{Reproducibility}

Code and data: \url{https://github.com/amitduabits/PRAHARI-P4-RetialQueues}

\begin{enumerate}
  \item \texttt{python instrument\_p4.py --k-values 1,2,4,8,12,16,24 --frames-per-k 200}
  \item \texttt{python plot\_frontier.py results/p4\_frontier.json}
\end{enumerate}

Expected runtime: 15–20 minutes on a modern laptop.

\section{Event Log Schema}

Each decode event:
\begin{lstlisting}
{
  "ts": "2026-09-03T11:20:30Z",
  "decode_request_id": "uuid",
  "K": 8,
  "latency_ms": 11.4,
  "cache_hit": true,
  "retried": false,
  "abandoned": false
}
\end{lstlisting}

\begin{thebibliography}{99}

\bibitem{erlang1917} Erlang, A.K. (1917). ``Solution of some problems in the theory of probabilities of significance in automatic telephone exchanges.'' Post Office Electrical Engineers' Journal.

\bibitem{retrial2010} Falin, G., Templeton, J.G. (2010). \emph{Retrial Queues}. Chapman \& Hall.

\bibitem{nvdec} NVIDIA. ``NVDEC: Video Decode GPU Engine.'' Developer documentation.

\bibitem{prahari} Dua, A. ``PRAHARI Platform.'' \url{https://github.com/amitduabits/PRAHARI}

\end{thebibliography}

\end{document}
